{
\selectlanguage{german}
\frontchapter{Zusammenfassung}

Diese Arbeit beschreibt das Design und die Implementation von Aesop, einer Taktik zur Automatisierung von Beweisen im interaktiven Theorembeweiser Lean~4.
Aesops Fundament ist eine Baumsuche im Stil von Tableaus, ähnlich derjenigen der auto-Taktiken von Rocq und Isabelle.
Nutzer können die Suche anpassen, indem sie Beweisregeln angeben, die Aesop dann anwendet.
Auf dieser Basis bauen drei wesentliche neue Funktionen auf.

Erstens unterstützt Aesop beliebige Suchstrategien und verwendet standardmäßig eine Best-First-Suchstrategie, wohingegen die meisten ähnlichen Taktiken auf Tiefensuche basieren.
Daraus folgt, dass die übliche, auf Backtracking basierende Methode, Metavariablen zu unterstützen, nicht mehr anwendbar ist.
Metavariablen sind \enquote*{Löcher} im Beweis und repräsentieren Subterme, die in den folgenden Beweisschritten auszufüllen sind.
Wir entwickeln eine neue Methode, die die Unabhängigkeit von Beweiszielen in unseren Suchbäumen auch dann bewahrt, wenn die Beweisziele Metavariablen enthalten, und dabei -- so unsere Vermutung -- vollständig bleibt.

Zweitens optimieren wir die Anwendung von Vorwärts-Regeln, die zusätzliche Fakten aus den lokalen Hypothesen ableiten.
Ähnliche Taktiken, die Vorwärts-Regeln unterstützen, verwenden eine recht naive, zustandslose Methode, um sie anzuwenden.
Dies führt zu erheblichen Redundanzen.
Wir adressieren diese mit einer Caching-Strategie, die strukturelle Ähnlichkeiten zwischen Beweiszielen ausnutzt.
Dies beschleunigt Aesop auf Beispielen, in denen viele Vorwärts-Regeln verwendet werden, um mehr als den Faktor 30.
Angewendet auf einen großen Benchmark, bei dem Vorwärts-Regeln weniger intensiv genutzt werden, erhalten wir im Durchschnitt eine Beschleunigung um den Faktor 1,24x.
Wie auch unsere Methode zur Unterstützung von Metavariablen ist der Algorithmus für Vorwärts-Regeln weitgehend unabhängig von der verwendeten Logik und könnte auch für nicht-typentheoretische Systeme verwendet werden.

Drittens implementieren wir eine Aesop-Funktion, die aus den angewendeten Regeln weitgehend idiomatische Taktik-Beweise generiert.
Das ermöglicht Nutzern, teure Aesop-Suchen durch die gefundenen Beweise zu ersetzen, oder auch Aesop als Präprozessor zu verwenden, der offensichtliche Beweisschritte automatisiert und dann die Kontrolle wieder an den Nutzer abgibt.
Um die generierten Beweise idiomatischer zu machen, implementieren wir einen Optimierer, der Taktiksequenzen in strukturierte Beweise transformiert und gewisse unidiomatische Konstruktionen eliminiert.
}
